\chapter{JSON, error boundaries, and middleware}
\chapterintro{A backend framework becomes useful when handlers can work with structured content, expected failures become consistent responses, and cross-cutting behavior can wrap every route.}

\begin{objectives}
\begin{itemize}
  \item Validate media type and UTF-8 before decoding JSON.
  \item Separate malformed content, invalid domain values, and unexpected failures.
  \item Compose middleware around a terminal router handler.
  \item Preserve a clear request-to-response data flow.
\end{itemize}
\end{objectives}

\section{JSON is one layer, not the wire protocol}

JSON parsing begins only after HTTP framing produces exact body bytes. The request
helper verifies that \code{Content-Type} names \code{application/json}, decodes
UTF-8, then calls the standard-library JSON parser. These failures have different
causes:

\begin{itemize}
  \item wrong or missing media type: 415;
  \item invalid UTF-8: 400;
  \item malformed JSON grammar: 400;
  \item valid JSON with wrong task fields: 422;
\end{itemize}

A JSON number, list, string, or null is syntactically valid. A Task endpoint
requires an object, so domain validation must still check the top-level type.

\section{A transitional application layer}

Stage 6 demonstrates JSON factories, a terminal dispatcher, and timing middleware
in one file:

\codefile{../code/stages/06_json_middleware/server.py}

The lambda defaults in \code{compose} are deliberate. Python closures capture
variables by reference; binding each current middleware and next handler as a
default prevents every wrapper from later referring to the final loop values.
The final \code{Application} uses a named wrapper with the same idea for clarity.

\section{Middleware is nested control flow}

Given middleware A and B, request flow is:

\begin{wiremessage}
A before
  B before
    route handler
  B after
A after
\end{wiremessage}

Each middleware receives a request and the next handler. It may modify the request
within documented bounds, return early, call the next handler once, inspect or
replace the response, or translate an expected exception. Logging, request IDs,
authentication, CORS, compression, and metrics are common examples, though not
all belong in this small project.

Order matters. An exception-recovery layer must surround code whose exceptions
it catches. A timing layer must begin before and finish after its measured work.
An authentication layer placed after a handler is not authentication.

\section{The final application boundary}

\codefile{../code/final/tinyhttp/application.py}

\code{Application.handle} is the safe public boundary. It rejects unsupported
methods, composes middleware, dispatches the route, verifies that handlers return
a Response, converts \code{HTTPError}, and catches unexpected exceptions.

Expected errors carry status, public detail, a stable machine code, optional
field errors, and a problem type URI. Unexpected exceptions are logged through
the standard \code{logging} package and become a generic 500 document. The catch
is broad only at this outer boundary; inner code should catch only errors it can
handle meaningfully.

\begin{securityboundary}
Never return \code{str(exception)} for an unexpected error. SQLite exceptions,
filesystem errors, and programmer mistakes can expose paths, schema details, or
secrets. Give clients a stable generic problem and give operators the traceback
through controlled logging.
\end{securityboundary}

\section{Stable errors are part of the API}

Clients should branch on status and stable problem type/code, not English prose.
The Task API uses a type such as
\code{urn:task-api:problem:validation-error} and includes an \code{errors} object
for individual fields. The prose remains useful to humans and may improve over
time without breaking client logic.

The \code{instance} identifies the request path associated with the occurrence.
In a larger system it could be a traceable occurrence URI, but it must not leak
private identifiers or internal topology.

\section{Observability without coupling}

The server access log records client address, method, target, status, response
bytes, and elapsed time. Application logs record unexpected failures. These are
different perspectives. A request ID middleware could connect them, but this
book avoids inventing a distributed tracing system.

Do not log complete request bodies by default. Bodies may contain credentials,
personal data, or large values. Log structure and identifiers deliberately,
with retention and access policy in mind.

\summarytitle

JSON decoding validates media type, encoding, and syntax before domain rules.
Expected failures become structured problem responses; unexpected failures stay
private. Middleware wraps the terminal dispatcher in a visible order and is the
framework's tool for reusable cross-cutting behavior.

\reviewtitle
\begin{enumerate}
  \item Why is a valid JSON list still invalid for task creation?
  \item What changes when two middleware functions exchange order?
  \item Why should clients not branch on an English error sentence?
  \item Which request data is dangerous to log by default?
\end{enumerate}

\begin{exercise}[Error taxonomy]
For twelve failures spanning HTTP syntax, JSON, task fields, missing tasks, SQL,
and programmer errors, choose the layer that owns detection, public status, stable
problem code, and whether a traceback is logged.
\end{exercise}

\begin{exercise}[Request ID middleware]
Design middleware that accepts a syntactically safe incoming request ID or
generates one, places it on the response, and makes it available to logs. Define
length and character limits and explain why arbitrary input must not be copied
into logs or fields.
\end{exercise}

